% Public text-only snapshot of the instructor's Module 12 VAE slide deck.
% Upstream: /Users/setup/Library/CloudStorage/Box-Box/Teaching/6050/LaTeX/Module 12/12.3S VAE.tex
% Upstream SHA-256: 499233c350074f16ab99d1f0bb50db0ad5ae3e1fc4967c93162ba63f942a95cd
% Local preamble dependencies, TikZ source, executable code, external assets,
% product-current claims, and claims without a sufficiently clear verification
% boundary are omitted. The book independently checks and derives every published
% claim, implementation, experiment, and figure.
\documentclass{article}
\usepackage{amsmath}
\usepackage{amssymb}
\usepackage[margin=1in]{geometry}
\usepackage{hyperref}

\title{Lecture 12.3: Generative AI and Variational Autoencoders}
\author{Heman Shakeri}
\date{}

\begin{document}
\maketitle

\section*{Public Snapshot Boundary}

This standalone snapshot preserves the instructor-authored VAE sequence and selected
equations. Diagrams, executable examples, product references, and categorical claims
about latent geometry have been removed. The retained material is a provenance record,
not a replacement for an independent derivation from primary sources.

\section{From Autoencoders to Generative Models}

A deterministic autoencoder maps an observation to a point and back:

\[
\mathbf{z}=f_{\phi}(\mathbf{x}),
\qquad
\hat{\mathbf{x}}=g_{\theta}(\mathbf{z}).
\]

Reconstruction training alone does not specify a probability distribution from which
new latent points should be sampled. A variational autoencoder instead introduces a
latent-variable model

\[
p_{\theta}(\mathbf{x},\mathbf{z})
=p(\mathbf{z})p_{\theta}(\mathbf{x}\mid\mathbf{z})
\]

and an approximate posterior $q_{\phi}(\mathbf{z}\mid\mathbf{x})$. A common choice is

\[
q_{\phi}(\mathbf{z}\mid\mathbf{x})
=\mathcal{N}\!\left(
\mathbf{z};
\boldsymbol{\mu}_{\phi}(\mathbf{x}),
\operatorname{diag}(\boldsymbol{\sigma}_{\phi}^2(\mathbf{x}))
\right),
\qquad
p(\mathbf{z})=\mathcal{N}(0,I).
\]

The lecture names Kingma and Welling's ``Auto-Encoding Variational Bayes'' as its
primary starting point: \url{https://arxiv.org/abs/1312.6114}.

% [The original architecture and latent-geometry diagrams are omitted.]

\section{The Evidence Lower Bound}

The marginal likelihood integrates over the latent variable:

\[
\log p_{\theta}(\mathbf{x})
=\log\int p_{\theta}(\mathbf{x},\mathbf{z})\,d\mathbf{z}.
\]

Introducing $q_{\phi}(\mathbf{z}\mid\mathbf{x})$ and applying Jensen's inequality
gives

\begin{align*}
\log p_{\theta}(\mathbf{x})
&=\log\int q_{\phi}(\mathbf{z}\mid\mathbf{x})
\frac{p_{\theta}(\mathbf{x},\mathbf{z})}
{q_{\phi}(\mathbf{z}\mid\mathbf{x})}\,d\mathbf{z}\\
&\geq
\mathbb{E}_{q_{\phi}(\mathbf{z}\mid\mathbf{x})}
\left[
\log p_{\theta}(\mathbf{x},\mathbf{z})
-\log q_{\phi}(\mathbf{z}\mid\mathbf{x})
\right]\\
&=
\mathbb{E}_{q_{\phi}(\mathbf{z}\mid\mathbf{x})}
[\log p_{\theta}(\mathbf{x}\mid\mathbf{z})]
-D_{\mathrm{KL}}\!\left(
q_{\phi}(\mathbf{z}\mid\mathbf{x})\,\|\,p(\mathbf{z})
\right).
\end{align*}

The first term rewards the decoder's conditional log likelihood. The second penalizes
departure of the approximate posterior from the prior. Their relative influence is an
optimization and modeling trade-off; neither term alone certifies useful semantics.

For a diagonal Gaussian posterior and standard-normal prior, the per-observation KL is

\[
D_{\mathrm{KL}}(q_{\phi}\|p)
=\frac{1}{2}\sum_{j=1}^{J}
\left(
\mu_j^2+\sigma_j^2-\log\sigma_j^2-1
\right).
\]

% [The original claim that the KL term makes every prior sample decode to a valid
% observation is omitted. Aggregate-posterior mismatch and decoder behavior remain
% empirical questions.]

\section{The Reparameterization Estimator}

For the Gaussian family above, sampling can be written as

\[
\boldsymbol{\epsilon}\sim\mathcal{N}(0,I),
\qquad
\mathbf{z}
=\boldsymbol{\mu}_{\phi}(\mathbf{x})
+\boldsymbol{\sigma}_{\phi}(\mathbf{x})\odot\boldsymbol{\epsilon}.
\]

The noise is sampled from a distribution independent of the encoder parameters. The
sample is then a differentiable function of those parameters and the sampled noise,
which supports pathwise gradient estimates.

% [The original reparameterization diagram, which closely follows a primary-paper
% figure, is omitted. Consult the cited primary literature for the original figure.]

\section{Weighted KL Objectives}

A $\beta$-weighted objective has the form

\[
\mathcal{L}_{\beta}
=\mathbb{E}_{q_{\phi}(\mathbf{z}\mid\mathbf{x})}
[\log p_{\theta}(\mathbf{x}\mid\mathbf{z})]
-\beta D_{\mathrm{KL}}\!\left(
q_{\phi}(\mathbf{z}\mid\mathbf{x})\,\|\,p(\mathbf{z})
\right).
\]

Changing $\beta$ changes the pressure toward the prior relative to reconstruction.
The original deck's categorical implication that this guarantees recovery of the
``true'' independent factors is omitted. Disentanglement depends on data, model,
inductive bias, supervision, objective, and evaluation protocol.
For the named objective, consult Higgins et al.,
\href{https://openreview.net/forum?id=Sy2fzU9gl}
{``beta-VAE: Learning Basic Visual Concepts with a Constrained Variational
Framework''}. For the identifiability boundary, consult Locatello et al.,
\href{https://proceedings.mlr.press/v97/locatello19a.html}
{``Challenging Common Assumptions in the Unsupervised Learning of Disentangled
Representations''}.

\section{Sampling and Interpolation}

After training, unconditional generation samples
$\mathbf{z}\sim p(\mathbf{z})$ and then samples or decodes from
$p_{\theta}(\mathbf{x}\mid\mathbf{z})$. Interpolating between two encoded points can
probe how the decoder behaves between them. Smooth-looking interpolation is a
diagnostic rather than proof that Euclidean movement has a prescribed semantic
meaning.

% [All executable implementation and visualization code is omitted.]

\section{Connection to Latent Diffusion}

A learned encoder--decoder can define a lower-dimensional state space for another
generative model. In latent diffusion, a diffusion process is trained on encoded
representations and generated representations are decoded afterward. The VAE and the
diffusion model retain distinct objectives; the connection does not make their latent
distributions or likelihood claims interchangeable.

% [Named-product examples and release-specific claims are omitted.]

\section{Limitations and Extensions}

The lecture closes with several recurring VAE concerns: observation models may favor
overly smooth reconstructions; a flexible decoder may ignore the latent variable
(posterior collapse); and the aggregate encoded distribution may not match the chosen
prior. Named extensions alter the prior, posterior family, objective, or latent
structure, but each change introduces assumptions that must be evaluated separately.

\section{Lecture Through-Line}

The lecture's sequence is: replace a deterministic code with an approximate posterior,
derive the ELBO, use reparameterization for gradient estimation, sample from a stated
prior, and examine the resulting trade-offs. The book may use this spine only after
rederiving the mathematics and testing any empirical claim independently.

\end{document}
